

\newpage

\section{Propagators testing (Phys. Rev. E 96, 063307 (2017))}

We want to solve the Schr\"odinger equation (\ref{eqn:tdse_2}) 
by treating some of these terms as inhomogeneous terms, namely

\begin{align}
 i\frac{\partial}{\partial t}\psi(x,t) & = \left[
	\left(-\frac{1}{2}\frac{\partial^2}{\partial x^2} 
           + \frac{1}{2}x^2\right)
           + \frac{1}{2}x
           +\frac{1}{2}\left(x +1\right)\right]\psi(x,t),
	  \hspace{9mm} \omega=1 \label{eqn:inhomo_fake} \\
              & = \Big( H_0 +V_{\rm ext} + h_I\Big)\psi(x,t).
\end{align}

\noindent The general theory for ordinary differential equations
consists in writing them as

\begin{equation}
   dy = dt\,f(t,y), \hspace{50pt} y(t_0)=y_0,
   \label{eqn:gen_ode}
\end{equation}


To this end, we assess several methods provided by the above-mentioned 
paper to numerically solve equations of the type 

\begin{equation}
   \frac{dy}{dt} = Ly + N(y,t),
\end{equation}

\noindent where $L$ is the linear part and $N(y,t)$ is the nonlinear part.
The equation we concern ourslevfes with is the 
time-dependent Kohn-Sham (TDKS) equation

\begin{align}
   i\frac{\partial}{\partial t}{\boldsymbol \rm \Psi}(\boldsymbol{r},t) 
& = \boldsymbol{\rm H}{\boldsymbol \rm \Psi}(\boldsymbol{r},t), 
\label{eqn:tdks} \\
& = \boldsymbol{\rm L}{\boldsymbol \rm \Psi}(\boldsymbol{r},t),
   +  \boldsymbol{\rm N}\left({\boldsymbol \rm \Psi}(\boldsymbol{r},t\right).
    \label{eqn:tdks}
\end{align}

\subsection{Analytical expressions}
\subsubsection{Time evolution operator approach in Schr\"odinger picture}

A formal solution to (\ref{eqn:tdks}) is 
\begin{equation}
   {\boldsymbol \rm \Psi}(\boldsymbol{r},t) = U(t,t_0)
   {\boldsymbol \rm \Psi}(\boldsymbol{r},t_0),
   \label{eqn:schrodinger_pic}
\end{equation}

\noindent where $U(t,t_0)$ is the evolution operator defined as

\begin{equation}
   U(t,t_0) = \mathcal{T}
     \exp\Bigg\lbrace 
     -i\int_{t_0}^{t}{\boldsymbol{\rm H}}(\boldsymbol{r},\tau)d\tau
\Bigg\rbrace,
\end{equation}

\noindent with $\mathcal{T}$ being the time-ordering operator.
In practice, the above expression is split into a product 
of evolution operators as

\begin{equation}
   U(t_n,t_0) = \prod_{k=1}^n U\left(t_k,t_{k-1}\right), \hspace{50pt}
         t_n = n\Delta t.
\end{equation}

\noindent For sufficiently small steps, we assume the Hamiltonian is kept 
constant between the two instants and write

\begin{equation}
   U(t_{n+1},t_n) = \exp\Big[
     -i{\boldsymbol{\rm H}}(\boldsymbol{r},t_n)\Delta t
\Big].
\end{equation}

\subsubsection{Implicit-explicit (IMEX) schemes}


The equation (\ref{eqn:gen_ode}) 
can be straightforwardly and conveniently solved 
by approximating it, for sufficiently small steps 
$\Delta t \longrightarrow 0$, to
  
\begin{equation}
   y_{n+1} = y_n + \Delta t\,f(t_n,y_n) + O(\Delta t^2), 
\hspace{50pt} y(t_0)=y_0.
   \label{eqn:explicit}
\end{equation}

\noindent This expression known as an \emph{explicit scheme}
interestingly consider that, for sufficiently small enough steps,
the function did not change much between $t_n$ and $t_{n+1}$. \\

An alternative to (\ref{eqn:explicit}) is the \emph{implicit scheme}
given by the following equation

\begin{equation}
   y_{n+1} = y_n + \Delta t\,f(t_{n+1},y_{n+1}) + O(\Delta t^2),
  \hspace{50pt} y(t_0)=y_0.
   \label{eqn:implicit}
\end{equation}

\noindent This last expression is preferable whenever 
the problem is said to be \emph{stiff}, as it is much stabler.



\subsubsection{Integrating factor (IF) method}

The idea here is to use an integrating factor to simplify the resolution 
of the equation.

\noindent As far as quantum mechanics is concerned,
it is about using the \emph{interaction picture} defined as the following
\begin{equation}
   \boldsymbol{\rm \Phi}(\boldsymbol{r},t) 
      = e^{i\boldsymbol{L}t}{\boldsymbol\rm\Psi} (\boldsymbol{r},t).
     \label{eqn:interac_pic}
\end{equation}

Inserting this into (\ref{eqn:tdks}) yields

\begin{equation}
   i\frac{\partial}{\partial t}\boldsymbol{\rm \Phi}(\boldsymbol{r},t) 
      = e^{i\boldsymbol{L}t}
    \boldsymbol{N}\left( e^{-i\boldsymbol{L}t}\boldsymbol{\rm\Phi},t\right).
     \label{eqn:base_if}
\end{equation}

Reversing back to $\boldsymbol{\rm \Psi}(\boldsymbol{r},t)$
form (\ref{eqn:interac_pic}) subsequently 
integrating from $t_n$ to $t_{n+1}$ yields

\begin{equation}
   \boldsymbol{\rm \Psi}(\boldsymbol{r},t_{n+1}) 
      = e^{-i\boldsymbol{L}\Delta t}
   \boldsymbol{\rm \Psi}(\boldsymbol{r},t_n) 
      -i e^{-i\boldsymbol{L}t_{n+1}}
      \int_{t_n}^{t_{n+1}} e^{i\boldsymbol{L}\tau}
    \boldsymbol{N}\left(  e^{-i\boldsymbol{L}\tau}
     \boldsymbol{\rm\Phi},\tau\right)d\tau.
     \label{eqn:if_step}
\end{equation}


\subsubsection{Exponential time differencing (ETD) method}

From (\ref{eqn:tdks}), we may express $\boldsymbol{N}(\boldsymbol{\Psi},t)$
as a sum of the linear part and the time-derivative. Subsequently
using the identity

\begin{equation}
   i\frac{\partial}{\partial t}
      e^{i\boldsymbol{L}t}\boldsymbol{\rm \Psi}(\boldsymbol{r},t) 
      = e^{i\boldsymbol{L}t}
    \Bigg[ 
    - \boldsymbol{L}(\boldsymbol{r},t)\boldsymbol{\rm \Psi}(\boldsymbol{r},t)
    + i\frac{\partial}{\partial t}
\boldsymbol{\Psi}(\boldsymbol{r},t)
    \Bigg],
\end{equation}

\noindent one can rewrite (\ref{eqn:tdks}) as 

\begin{equation}
   i\frac{\partial}{\partial t}\boldsymbol{\rm \Phi}(\boldsymbol{r},t) 
      = e^{i\boldsymbol{L}t}
    \boldsymbol{N}\left( \boldsymbol{\rm\Psi},t\right),
     \label{eqn:base_etd}
\end{equation}

\noindent Integrating from $t_n$ to $t_{n+1}$ yields

\begin{equation}
   \boldsymbol{\rm \Psi}(\boldsymbol{r},t_{n+1}) 
      = e^{-i\boldsymbol{L}\Delta t}
   \boldsymbol{\rm \Psi}(\boldsymbol{r},t_n) 
      -i e^{-i\boldsymbol{L}t_{n+1}}
      \int_{t_n}^{t_{n+1}} e^{i\boldsymbol{L}\tau}
    \boldsymbol{N}\left( \boldsymbol{\rm\Psi},\tau\right)d\tau.
     \label{eqn:etd_step}
\end{equation}

\noindent The expressions (\ref{eqn:base_etd}) and (\ref{eqn:etd_step})
are formally identical to (\ref{eqn:base_if}), (\ref{eqn:if_step}),
safe the variable change in the treatment of the nonlinear term.
In other words, the IF method consists in propagating the solution 
in the interaction picture whereas the ETD method does so in
the Schr\"odinger picture.

\vspace{20pt}

From this last remark, there is an argument to make about the point 
of propagating the solution by the means of (\ref{eqn:if_step}) or
(\ref{eqn:etd_step}). Indeed, one can just propagate 
the solution in the interaction picture all the way through the last 
integration step 

\begin{equation}
   \boldsymbol{\rm \Phi}(\boldsymbol{r},t_{n+1}) 
      = \boldsymbol{\rm \Phi}(\boldsymbol{r},t_n) 
      -i \int_{t_n}^{t_{n+1}} e^{i\boldsymbol{L}\tau}
    \boldsymbol{N}\left(  e^{-i\boldsymbol{L}\tau}
     \boldsymbol{\rm\Phi},\tau\right)d\tau,
\end{equation}

\noindent and then, and then only, obtain the solution
$\boldsymbol{\rm \Psi}(\boldsymbol{r},t_f)$ from 
$\boldsymbol{\rm \Phi}(\boldsymbol{r},t_f)$  by making use of 
(\ref{eqn:interac_pic}).

\newpage

\subsection{Solutions to the Time Dependent Density Functional Theory}

In this part, we address some specific technics to solve 
the TDKS equation, 
either in the Schr\"odinger picture (\ref{eqn:schrodinger_pic})
or in the interaction picture (\ref{eqn:interac_pic}).

\subsubsection{Taylor time propagation}

From the expression (\ref{eqn:schrodinger_pic}),
a simple Taylor expansion yields

\begin{equation}
\boldsymbol{\rm \Psi}(\boldsymbol{r},t_{n+1})
   = \sum_{k=0}^\infty \frac{1}{k!}
      \Bigg( -i\Delta t\boldsymbol{\rm H}(t_n)\Bigg)^k
     \boldsymbol{\rm \Psi}(\boldsymbol{r},t_n)
     \label{eqn:taylor}
\end{equation}

\noindent The developpment is usually truncated after $n$ terms,
which in turns gives an  $n$th-order approximation.

\begin{equation}
\boldsymbol{\rm \Psi}(\boldsymbol{r},t_{n+1})
   = \sum_{k=0}^n \frac{1}{k!}
      \Bigg( -i\Delta t\boldsymbol{\rm H}(t_n)\Bigg)^k
     \boldsymbol{\rm \Psi}(\boldsymbol{r},t_n) + O(\Delta t^{n+1}).
\end{equation}

\noindent Most propagation schemes seems to be essentially just that 
in some way or another.

\subsubsection{Split-operator method}
\subsubsection{Crank-Nicholson method}

By making use of both the explicit and implicit schemes and 
averaging them, we achieve the Crank Nicholson propagation scheme 
given by the following formula

\begin{\eq}
  y_{n+1} = y_n + \frac{\Delta t}{2}
  \Bigg(f(t_{n+1},y_{n+1}) + f(t_n,y_n)\Bigg) + O(\Delta t^3)
  \label{eqn:crank_nich}
\end{\eq}

\noindent The application of this formula to the 
problem (\ref{eqn:tdks}) gives 

\begin{align}
   \boldsymbol{\rm \Psi}_{n+1}  =
   \boldsymbol{\rm \Psi}_n
   - \frac{i\Delta t}{2}\Bigg( \boldsymbol{H}_{n+1}\boldsymbol{\rm \Psi}_{n+1}
   + \boldsymbol{H}_{n}\boldsymbol{\rm \Psi}_{n}\Bigg), &
    \hspace{50pt} \boldsymbol{H}_i \equiv \boldsymbol{H}(\boldsymbol{r},t_i)\\
   \Bigg\lbrace\mathds{1} +  \frac{i\Delta t}{2}
   \Big( \boldsymbol{L}_{n+1} + \boldsymbol{N}_{n+1}\Big)\Bigg\rbrace
   \boldsymbol{\rm \Psi}_{n+1}
   & = \Bigg\lbrace\mathds{1} +  \frac{i\Delta t}{2}
       \Big( \boldsymbol{L}_{n} - \boldsymbol{N}_{n}\Big)\Bigg\rbrace
       \boldsymbol{\rm \Psi}_{n}
\end{align}

\subsubsection{Second-order IMEX method}
\subsubsection{Integrating factor method with explicit multistep}

This method is not derived from (\ref{eqn:explicit}) but rather 
from the \emph{Adam-Bashfort} formula, which reads

\begin{\eq}
   y_{n+1} = y_n 
      + \frac{\Delta t}{2}\Bigg( 3f(y_n,t_n) - f(y_{n-1},t_{n-1})\Bigg) + O(\Delta t^3) ,
     \hspace{50pt} \Delta t = t_{n+1} - t_{n}.
\end{\eq}

The application of this formula to (\ref{eqn:interac_pic}) gives 

\begin{align}
   \boldsymbol{\rm \Phi}_{n+1} = \boldsymbol{\rm \Phi}_{n} 
      + \frac{i\Delta t}{2}
\Bigg[ -3e^{i\boldsymbol{L}t_n}\boldsymbol{N}\left( e^{-i\boldsymbol{L}t_n}
       \boldsymbol{\Phi}_n,t_n\right)
+ e^{i\boldsymbol{L}t_{n-1}}\boldsymbol{N}\left( e^{-i\boldsymbol{L}t_{n-1}}
       \boldsymbol{\Phi}_{n-1},t_{n-1}\right)\Bigg], \\
\boldsymbol{\rm \Psi}_{n+1} = 
      e^{-i\boldsymbol{L}\Delta t}\boldsymbol{\rm \Psi}_{n} 
 -\frac{3i\Delta t}{2} e^{-i\boldsymbol{L}\Delta t}
       \boldsymbol{N}\left(\boldsymbol{\Psi}_n,t_n\right)
+ \frac{i\Delta t}{2} e^{-2i\boldsymbol{L}\Delta t}\boldsymbol{N}
\left(\boldsymbol{\Psi}_{n-1},t_{n-1}\right).
\end{align}

\subsubsection{Integrating factor with Runge-Kutta}
\emph{Runge Kutta} methods propagate the solution over an interval
by combining the information from several Euler-style steps (each involving
one evaluation of the right-hand function), 
and then using the information obtained to matrch a Taylor series 
expansion up to some order. \\

%In other words Runge-Kutta method and Euler method are fundamentally the same,
%as can been seen by the following development. 

\begin{itemize}
\item \textbf{Euler method (RK1)}

The formula for the Euler method is given by
\begin{align}
   k_1 & =  f(t_n,y_n), \label{eqn:rk1_step}\\
   y_{n+1} & = y_n + \Delta t\, k_1 + O(\Delta t^2) \label{eqn:rk1}.
\end{align}
The formula for the explicit Euler method (\ref{eqn:explicit})
advances the solution one step forward in time. The method is described
to be $O(\Delta t^2)$ or $o(\Delta t)$, thus of order 1.
Though imprecise and unstable, Euler method may be used
to obtain a partial solution that is then used in more advanced 
propagation schemes.

\begin{align}
        k_1 & = -i e^{-i\boldsymbol{L}\Delta t}
        \boldsymbol{N}(\boldsymbol{\rm\Psi}_{n},t_n), \\
   \boldsymbol{\rm\Psi}_{n+1} 
      & = e^{-i\boldsymbol{L}\Delta t}\boldsymbol{\rm\Psi}_{n}
      +  \Delta t\, k_1.
\end{align}

\hspace{20pt}

\item \textbf{Second-order Runge-Kutta methods}

Using (\ref{eqn:rk1_step}) as a trial step to some point of the interval 
and subsequently using the value of $t$ and $y$ at this point, we may 
compute the real step as

\begin{align}
   k_1 & =  f(y_n,t_n), \label{eqn:rk1_step } \\
   k_2 & =  f\left(t_n + \alpha\Delta t,y_n+\alpha\Delta t\,k_1\right), 
      \label{eqn:rk2_step }\\
   y_{n+1} & = y_n + \Delta t\Big((1-\beta) k_1 
              +  \beta k_2 \Big) + O(\Delta t^3) \label{eqn:rk2},
\end{align}
\noindent where $\beta=(1-\alpha)$ and $\alpha$ is 
some non-zero real parameter.
For example $\alpha=1/2$ is the second-order Runge's method 
also known as the \emph{midpoint rule}. Other popular choices
are $\alpha=2/3$  which is Ralston's method
and $\alpha=1$ which is Heun's method. \\

Using (\ref{eqn:rk2}) on the relation (\ref{eqn:interac_pic}) yields
\begin{align}
        k_1 & = -i e^{-i\boldsymbol{L}\Delta t}
        \boldsymbol{N}(\boldsymbol{\rm\Psi}_{n},t_n), \\
        k_2 & = -i e^{-i\alpha\boldsymbol{L}\Delta t}
        \boldsymbol{N}(\boldsymbol{\rm\Psi}_{n} + \alpha\Delta t\,k_1,
             t_n + \alpha\Delta t), \\
   \boldsymbol{\rm\Psi}_{n+1} 
      & = e^{-i\boldsymbol{L}\Delta t}\boldsymbol{\rm\Psi}_{n}
      +  \Delta t\Big( (1-\beta)k_1 + \beta k_2\Big).
\end{align}

\hspace{20pt}
\item \textbf{Third-order Runge-Kutta methods}

Following the second-order Runge-Kutta methods 
are the third-order methods described by the equations

\begin{align}
   k_1 & =  f(y_n,t_n), \label{eqn:rk1_step } \\
   k_2 & =  f\left(t_n + c_2\Delta t,y_n + a_{21}\Delta t\,k_1\right), \\
   k_3 & =  f\big(t_n + c_3\Delta t, 
        y_n + \Delta t\left( a_{31}k_1 +a_{32}k_2\right)\big), 
      \label{eqn:rk3_step }\\
   y_{n+1} & = y_n + \Delta t\Big( b_1k_1 
              +  b_2 k_2 + b_3 k_3\Big) + O(\Delta t^4) \label{eqn:rk3},
\end{align}

\noindent where the coefficients $c_i$, $a_{ij}$ and $b_j$
are given by the \emph{Butcher tableau}. 

\begin{table}%[!h]
\begin{tabular}{l|c}
$\boldsymbol{\rm c}$      &   $\boldsymbol{A}$                   \\
\hline
\rule{0pt}{3ex}           &  $\boldsymbol{\rm b^T}$  
\end{tabular}
\end{table}


\hspace{20pt}
\item \textbf{Fourth-order Runge-Kutta methods}

Following are the fourth-order Runge-Kutta methods described by the equations

\begin{align}
   k_1 & =  f(y_n,t_n), \label{eqn:rk1_step } \\
   k_2 & =  f\left(t_n + c_2\Delta t,y_n + a_{21}\Delta t\,k_1\right), \\
   k_3 & =  f\big(t_n + c_3\Delta t, 
        y_n + \Delta t\left( a_{31}k_1 +a_{32}k_2\right)\big), \\
   k_4 & =  f\big(t_n + c_4\Delta t, 
        y_n + \Delta t\left( a_{41}k_1 +a_{42}k_2 + a_{43}k_3\right)\big), 
      \label{eqn:rk4_step }\\
   y_{n+1} & = y_n + \Delta t\Big( b_1k_1 
              +  b_2 k_2 + b_3 k_3 + b_4 k_4\Big) 
+ O(\Delta t^5) \label{eqn:rk4},
\end{align}

\noindent where the coefficients $c_i$, $a_{ij}$ and $b_j$
are, again, given by the \emph{Butcher tableau}. 

There exist two popular fourth-order Runge-Kutta methods,
only differing by the coefficients of the Butcher tableau. 
\end{itemize}

\begin{table}[!h]
\centering
%\begin{tabular}{cc}%
%\begin{tabular}{|@{\hspace{2em}}l@{}|l@{\qquad}|}
%\begin{tabular}{@{\hspace{2em}}c@{}@{\qquad}c}
\begin{tabular}{@{\hspace{2em}}c@{}@{\hspace{10em}}c}
\begin{tabular}{l|cccc}
$0$      &  $0$    & $0$   &  $0$   & $0$              \\
$1/2$    &  $1/2$  & $0$   &  $0$   & $0$              \\
$1/2$    &  $0$    & $1/2$ &  $0$   & $0$              \\
$1$      &  $0$    & $0$   &  $1$   & $0$              \\
\hline
         &  $1/6$  & $1/3$ &  $1/3$ & $1/6$              
\end{tabular} &
\begin{tabular}{l|cccc}
$0$      &  $0$    & $0$   &  $0$   & $0$              \\
$1/3$    &  $1/3$  & $0$   &  $0$   & $0$              \\
$2/3$    &  $-1/3$ & $1$   &  $0$   & $0$              \\
$1$      &  $1$    & $-1$  &  $1$   & $0$              \\
\hline
         &  $1/8$  & $3/8$ &  $3/8$ & $1/8$              
\end{tabular}
\end{tabular}
\end{table}

